\documentclass[a4paper,12pt]{article}
\usepackage{amsmath,amssymb,amsthm}
\usepackage{xcolor}
\usepackage{mdframed}
\usepackage{thmtools}

% Define color palette matching the image
\definecolor{thmorange}{HTML}{FFEAD5} % Theorem background
\definecolor{exred}{HTML}{FFDCDD}     % Example background
\definecolor{lemblue}{HTML}{E2E0FF}   % Lemma background
\definecolor{corpink}{HTML}{FFF0F2}   % Corollary background
\definecolor{defgray}{HTML}{E2E2E2}   % Definition background

% Generic style generator for solid background boxes
\mdfdefinestyle{solidbox}{
	hidealllines=true,
	innerleftmargin=12pt,
	innerrightmargin=12pt,
	innertopmargin=10pt,
	innerbottommargin=10pt,
	skipabove=12pt,
	skipbelow=12pt
}

% Define theorem styles using thmtools + mdframed
\declaretheoremstyle[
headfont=\bfseries,
bodyfont=\itshape,
mdframed={style=solidbox, backgroundcolor=thmorange}
]{thmstyle}

\declaretheoremstyle[
headfont=\bfseries,
bodyfont=\normalfont,
mdframed={style=solidbox, backgroundcolor=exred}
]{exstyle}

\declaretheoremstyle[
headfont=\bfseries,
bodyfont=\itshape,
mdframed={style=solidbox, backgroundcolor=lemblue}
]{lemstyle}

\declaretheoremstyle[
headfont=\bfseries,
bodyfont=\itshape,
mdframed={style=solidbox, backgroundcolor=corpink}
]{corstyle}

\declaretheoremstyle[
headfont=\bfseries,
bodyfont=\normalfont,
mdframed={style=solidbox, backgroundcolor=defgray}
]{defstyle}

% Declare theorem environments
\declaretheorem[style=thmstyle, name=Theorem, numberwithin=section]{theorem}
\declaretheorem[style=exstyle, name=Example, numberwithin=section]{example}
\declaretheorem[style=lemstyle, name=Lemma, numberwithin=section]{lemma}
\declaretheorem[style=corstyle, name=Corollary, numberwithin=section]{corollary}
\declaretheorem[style=defstyle, name=Definition, numberwithin=section]{definition}

\begin{document}
	
	\section{Introduction}
	
	\begin{theorem}
		If $V \subseteq U$, and if $F$ is a set of binary functions defined on $U$, then there exists a unique $\operatorname{clos}_F(V)$ and it is closed under $F$.
	\end{theorem}
	
	\begin{example}
		Let $V = \{\{1,2\}, \{2,3\}\}$, and $F$ be the set containing the set-union and set-intersection operations, denoted $F = \{\cup, \cap\}$. Then $\operatorname{clos}_F(V) = \{\emptyset, \{1,2\}, \{2\}, \{2,3\}, \{1,2,3\}\}$, because if we take $\alpha, \beta \in \{\emptyset, \{1,2\}, \{2\}, \{2,3\}, \{1,2,3\}\}$ then both $\alpha \cup \beta$ and $\alpha \cap \beta$ are also therein.
	\end{example}
	
	\begin{lemma}
		If $V \subseteq U$, and $F$ is the set of three primitive set operations union, intersection, and relative complement, ($F = \{\cup, \cap, \setminus\}$) then we denote $\operatorname{clos}_F(V)$ simply by $\sigma(V)$ and call it the Sigma algebra of $V$. Moreover, each element of $\sigma(V)$ is called a Boolean combination of elements of $V$.
	\end{lemma}
	
	\section{Conclusion}
	
	\begin{corollary}
		If $V \subseteq U$, $\sigma(V)$ exists and is unique.
	\end{corollary}
	
	\begin{definition}
		Let $V = \{\{1,2\}, \{2,3\}\}$ as in Example 1.1. $\sigma(V) = \{\emptyset, \{1\}, \{2\}, \{3\}, \{1,2\}, \{1,3\}, \{2,3\}, \{1,2,3\}\}$. $\sigma(V) = \rho(\{1,2,3\})$.
	\end{definition}
	
\end{document}